% Generated draft: requires mathematical and source review.
\documentclass[11pt]{article}
\usepackage[margin=25mm]{geometry}
\usepackage{fontspec}
\setmainfont{Latin Modern Roman}
\usepackage{amsmath,amssymb,mathtools,unicode-math}
\setmathfont{Latin Modern Math}
\usepackage{enumitem,needspace,xurl,hyperref}
\setlength{\parindent}{0pt}
\setlength{\parskip}{5pt}
\setlength{\emergencystretch}{3em}
\newcommand{\N}{\mathbb{N}}
\newcommand{\Z}{\mathbb{Z}}
\newcommand{\Q}{\mathbb{Q}}
\newcommand{\R}{\mathbb{R}}
\newcommand{\C}{\mathbb{C}}
\newcommand{\F}{\mathbb{F}}
\newcommand{\A}{\mathbb{A}}
\newcommand{\PP}{\mathbb{P}}
\newcommand{\E}{\mathbb{E}}
\newcommand{\eps}{\varepsilon}
\newcommand{\dd}{\,\mathrm{d}}
\begin{document}
\section*{UnsolvedMath ID: 30004888}
\textbf{Source problem code:} catalogue entry
\subsection{Definitions and mathematical statement}
Let \(L\) be a link and let \(H_y(L)\) denote its y-ified Khovanov--Rozansky homology in the construction intended by the supplied source. The source describes operators \(x_i,y_i\), a family \(F_k\) for \(k=1,2,\ldots\), and operators \(f,h\) associated with an \(\mathfrak{sl}_2\)-action extending \(F_2\). The meaning and range of the index \(i\), the coefficient system, the gradings, and the precise homology construction are not specified in the supplied record.

For each such operator action, write
\[
\mathcal A_L=\bigl\langle x_i,y_i,F_k,f,h\bigr\rangle
\]
for the associative algebra generated by these operators under addition, scalar multiplication in the intended coefficient system, and composition. A relation means a finite linear combination of words in these generators whose induced operator on \(H_y(L)\) is zero. Once the coefficient system and generator index sets are fixed, a presentation amounts to determining the kernel of the evaluation homomorphism from the corresponding free associative algebra to \(\mathcal A_L\).

\textbf{Problem.} Determine the full algebra of relations among \(x_i,y_i,F_k,f,h\) acting on y-ified Khovanov--Rozansky homology. The source does not specify whether the desired relations are all relations for each individual link \(L\), the relations holding simultaneously for every link, or a presentation of a separately defined algebra acting on every such homology. These targets need not coincide, so this distinction remains unresolved.

The supplied dated assessment lists the relations
\[
[F_k,F_m]=0,\qquad [F_k,x_i]=0,\qquad
[F_k,y_i]=k x_i^{k-1}
\]
for all applicable \(i\) and positive integers \(k,m\), and states that \(F_2\) extends to an \(\mathfrak{sl}_2\)-action. Here the displayed brackets reproduce the supplied commutator notation; under the ordinary convention, \([A,B]=AB-BA\). The record does not settle whether graded commutators are intended, nor does it specify the normalization identifying \(F_2,f,h\) with an \(\mathfrak{sl}_2\)-triple. Consequently, no additional bracket relations involving \(f\) or \(h\) are imposed in this draft.

As background, the extracted source states that any complex of the form \(\mathrm{HH}^{i}(F_\beta)\) admits an action of a differential graded algebra denoted \(\mathrm{CA}\), and that the resulting \(\mathrm{CA}\)-module is a link invariant. It describes passage to y-ified homology as a certain form of Koszul duality and states that the y-ified Khovanov--Rozansky homology of any link admits commuting operators \(F_1,F_2,\ldots\). The definitions of \(\mathrm{HH}^{i}(F_\beta)\), \(\beta\), and \(\mathrm{CA}\), and the relevant equivalence notion for the module invariant, are absent. The accompanying extracted formula for \(F_k\) contains displaced summation and derivative symbols and an undefined term \(u_k\); it does not provide a reliable operator definition. These background assertions and the listed relations are supplied source claims, not independently verified results in this draft.
\subsection{Short English statement}
Determine all relations among the operators \(x_i,y_i,F_k,f,h\) acting on y-ified Khovanov--Rozansky homology. The supplied record lists some commutator relations and an \(\mathfrak{sl}_2\)-extension of \(F_2\), but leaves the coefficient system, operator conventions, and intended meaning of a complete presentation unspecified. It also leaves unclear whether the target concerns each link separately or relations valid for every link.
\subsection{Sources}
\begin{itemize}
\item Supplied local UnsolvedMath registry record, permanent ID \(30004888\), title: \emph{Operator Relations on Yified Khovanov--Rozansky Homology}; exact source problem code: \texttt{catalogue entry}. No individual proposer is supplied; the recorded proposed year is 2021.
\item Supplied report citation: \emph{Enveloping Algebras and Geometric Representation Theory} (2022), Oberwolfach Reports, pp.\ 2827--2891, DOI \texttt{10.4171/owr/2021/52}. \url{https://doi.org/10.4171/owr/2021/52}. The supplied citation uses 2022, while the DOI and registry year information refer to 2021; these details are retained without reconciliation.
\item Supplied literature reference labeled Gorsky--Hogancamp--Mellit: \url{https://arxiv.org/abs/2103.01212}. The record supplies no title or pinpoint reference for the asserted operator relations.
\item Supplied literature reference labeled OWR report: \url{https://doi.org/10.4171/owr/2021/52}.
\item Literature assessment date: \texttt{2026-08-22}. This is supplied dated curation, not an independently verified assessment. It labels the problem partially solved, lists the displayed commutator relations and the \(\mathfrak{sl}_2\)-extension, and states that the complete algebra of relations remains open. The background explicitly calls this dated literature triage derived from supplied OWR curation and requests independent verification of time-sensitive status and cited claims.
\item The registry labels the statement \texttt{corrected\_verified} and reports that the official Oberwolfach report was checked directly and the clean statement restored definitions and removed extraction artifacts or source-dependent wording. This is a supplied verification claim; no source visit or independent verification is claimed here.
\end{itemize}
\end{document}
